% =====================================================================
% P1 Appendix B.1 --- Protocol details.
%
% Source of truth: paper/main-draft.md Section 6.2 (pre-M5 text) and
% results/generation_summary.json.  This file is a fragment: it is meant
% to be \input under \appendix after Appendix A.
%
% Counts and hyperparameters here must match Section 6.1 and the
% generation summary; do not change one without the other.
% =====================================================================

\section{Protocol details and additional results}\label{app:b}

\subsection{Protocol details}\label{app:b1}

\paragraph{L0 (row level).}
For each held-out regime $r$, train on a chronological 70\% of the rows from the
other four regimes, calibrate on the remaining 30\%, and test on all rows of $r$.
The iid baseline is a 70/15/15 random split that ignores regime. The compared
rules are a pooled radius, Mondrian radii per calibration regime with max
fallback, a similarity-class fallback ($\{\text{OPEN},\text{EXTENDED}\}$ versus
$\{\text{WEEKEND},\text{HOLIDAY},\text{HALT}\}$), and a strategy-conditional
radius.

\paragraph{L1 (trajectory level).}
Calibration and evaluation use independent path sets: 128 seeds $\times$ 4
windows each, for 512 calibration and 512 evaluation paths, with $R=32$
simulations per path. We use split conformal at $\alpha=0.1$,
raw/studentized/CQR scores, and a Gaussian plug-in normal baseline. The
protocols are within-window and within-regime positive controls and
leave-one-window-out/leave-one-regime-out. The paired bootstrap resamples test
paths 2{,}000 times, conditioning on the fitted calibration quantiles.

\paragraph{L2 (branch level).}
The event study uses the 128 offline paths (32 seeds $\times$ 4 windows), one
freeze-to-OPEN event per path at $\tau=2$ (main) and $\tau=81$ (transfer), with
five-fold out-of-fold predictions by path. \texttt{post\_rebalance} is used only
to evaluate coverage; the branch-union threshold
$\max(\hat q_{\mathrm{hold}},\hat q_{\mathrm{rebalance}})$ is computed on the
calibration side. Pre-decision Mondrian stratifies on \texttt{onset\_action},
\texttt{onset\_distance}, and \texttt{onset\_side}.

\paragraph{L3 (observation map).}
The observable-only estimate maximizes a state-space Gaussian-sum likelihood
with a two-component jump mixture and a trimmed likelihood that masks $|z|>3$.
Uncertainty is a bootstrap and a numerical Hessian. The five coverage specifications are the true-$S$ diagnostic (G0) plus four
observable-only re-runs that vary the trimming threshold ($z=3$ main, 2.5, 4)
and the volatility source (oracle or estimated from OPEN). An unmasked MLE is
used only in the parameter-specification curve and has no coverage re-run.

\paragraph{Second synthetic sensor-POMDP testbed.}
The latent state and action-dependent transition are
\[
x_{t+1}=0.97x_t+0.35a_t+0.20w_t,\qquad w_t\sim N(0,1).
\]
Steps 16--31 are degraded, steps 32--39 are outages, and all others are normal.
Normal observations are $y_t=x_t+\varepsilon_t$ with
$\varepsilon_t\sim N(0,0.10^2)$. Degraded observations are $y_t=x_t+b_t$ with
$b_t=0.85b_{t-1}+\nu_t$, $\nu_t\sim N(0,0.35^2)$. Outage observations are
$y_t=y_{31}+\eta_t$, $\eta_t=0.90\eta_{t-1}+\zeta_t$,
$\zeta_t\sim N(0,0.05^2)$. The logging action in the eight-step burn-in is
$\operatorname{clip}(-0.15y_t+e_t,-1,1)$ with $e_t\sim N(0,0.30^2)$; the target
action is $\operatorname{clip}(-0.70m_t,-1,1)$, where $m_t$ is the identified
belief mean. Truth and both predictive models use this target-action rule and
identified-belief filter, but action realizations follow their simulated
observation histories. The per-step reward is
$-(x_{t+1}^2+0.05a_t^2)$. The reported error is the mean absolute difference
between the 64-rollout predictive mean and realized target return $G$, not a
value-bias estimate; it is a descriptive point estimate without a paired
uncertainty interval. Both predictors use the known control coefficient $B=0.35$.
The identified-transition predictor also uses known $A$ and $\sigma$ and an
augmented Kalman belief in normal/degraded regimes, propagating the state belief
without an outage measurement update. The fitted-transition comparator uses the
same identified-belief target-action rule and filter; only its autonomous rollout
AR(1) coefficient and residual scale are estimated from burn-in observations/actions.
It is not an alternative observation-filter
or policy comparison. Calibration and evaluation use disjoint sets of 128 paths,
with 64 target rollouts per model and path. The intervention branch is any target
action with magnitude at least 0.65; pre-decision stratification uses three groups
from the sign and magnitude of the last burn-in observation. Random streams for
truth and estimation are separate, and the latent start is used only for the truth rollout.

\paragraph{I5 stress variant.}
A separate environment uses 384-step windows, 128 seeds $\times$ 4 windows $=$
512 paths, and 1{,}536 double-endpoint events at $\tau\in\{2,65,81\}$ (512 each),
with five-fold out-of-fold predictions by seed mod 5. An unfreeze jump
$J_e\sim N(0,g(\tau_e,\sigma_e)^2)$ with
$g=\gamma\sigma\sqrt{1-e^{-\tau/\tau_0}}$ and $\tau_0=40$ is injected only on
freeze-to-OPEN transitions. The $\gamma=4$ environment is compared with a
$\gamma=0$ placebo; twin paths share the main RNG stream. The cross-event
estimator fits $\operatorname{Var}(r_e)=v_{0e}+\theta w_e$ with
$w_e=\sigma^2(1-e^{-\tau_e/\tau_0})$ and $\theta=\gamma^2$.

\subsection{Additional tables moved from the main text}
\label{app:b-tables}

These tables carry the full panels behind the summaries in
Sections~\ref*{sec:identification}--\ref*{sec:results}. The moved tables are
B1--B9: B1--B4 are the three-term decomposition, bridge, semi-synthetic, and misspecification
tables, B5 is the L0/L1 coverage panel, B6 is the related-work positioning
table, B7 is the second synthetic sensor-POMDP testbed, B8 answers Q1--Q3, and B9 is
the real-path usage check.

\renewcommand{\thetable}{B\arabic{table}}
\renewcommand{\theHtable}{B\arabic{table}}  % unique hyperref anchors: avoid name{table.N} collisions with Appendix A/C/D tables
\setcounter{table}{0}
\begin{table*}[!htbp]
\centering
\scriptsize
\caption{Three-term decomposition, reported separately}
\label{tab:b1}
\par\noindent\begin{minipage}{\linewidth}\raggedright{\scriptsize The decomposition is not a theorem. The terms are along observed histories, not a uniform sup over beliefs; the third term is a diagnostic refinement for observation-type policies, not an independent bound term. The phases are anchor = last informative step before the freeze, freeze = pre+1..post-1, and post = first OPEN after the freeze.}\end{minipage}\par

\par\textbf{Panel A. $W_1$ estimation term (identified posterior vs naive delta; oracle uses latent $S_t$)}

\begin{adjustbox}{max width=\linewidth}
\begin{tabular}{@{}llllllll@{}}
\toprule
$\tau$ (h) & Phase & $n$ & $W_1$ mean & $W_1$ SD & $W_1$ p90 & Oracle $W_1$ mean & Ratio \\
\midrule
$\tau=2$ & anchor & 128 & 3.96e-04 & 1.63e-19 & 3.96e-04 & 4.04e-04 & 0.980 \\
$\tau=2$ & freeze & 128 & 0.0076 & 0.0028 & 0.0105 & 0.0077 & 0.987 \\
$\tau=2$ & post & 128 & 3.96e-04 & 5.44e-07 & 3.97e-04 & 4.10e-04 & 0.966 \\
$\tau=81$ & anchor & 128 & 0.0027 & 7.56e-05 & 0.0028 & 0.0028 & 0.953 \\
$\tau=81$ & freeze & 128 & 0.0407 & 0.0160 & 0.0567 & 0.0406 & 1.003 \\
$\tau=81$ & post & 128 & 3.96e-04 & 2.70e-08 & 3.96e-04 & 4.09e-04 & 0.968 \\
\bottomrule
\end{tabular}
\end{adjustbox}
\vspace{0.5em}
\par\noindent\begin{minipage}{\linewidth}\raggedright{\scriptsize Notes. $W_1 = \mathbb{E}|\log S_t - \log P_t|$ under the identified posterior, i.e., $W_1$ between the identified posterior and the delta at $\log P_t$ (the naive treats $P_t$ as $S_t$); the oracle replaces $P_t$ by the simulator's latent $S_t$ for validation only. Ratio = identified mean / oracle mean. $n$ is the number of episodes (128 per cell); the $\tau=2$ anchor cell is numerically exact (SD $<10^{-15}$); the $\tau=81$ anchor SD is $7.6\times10^{-5}$. The freeze term grows with the freeze duration: 0.0076 at $\tau=2$ and 0.0407 at $\tau=81$.}\end{minipage}\par

\par\textbf{Panel B. Predictive-distribution TV information term}

\begin{adjustbox}{max width=\linewidth}
\begin{tabular}{@{}llllllll@{}}
\toprule
$\tau$ (h) & Phase & $n$ & TV mean & TV SD & TV p90 & TV max & Step frac TV$\ge$0.99 \\
\midrule
$\tau=2$ & anchor & 128 & 1.000 & 0.000 & 1.000 & 1.000 & 1.00 \\
$\tau=2$ & freeze & 128 & 0.646 & 0.039 & 0.686 & 0.721 & 0.50 \\
$\tau=81$ & anchor & 128 & 0.739 & 0.261 & 1.000 & 1.000 & 0.27 \\
$\tau=81$ & freeze & 128 & 0.766 & 0.066 & 0.838 & 0.911 & 0.00 \\
\bottomrule
\end{tabular}
\end{adjustbox}
\vspace{0.5em}
\par\noindent\begin{minipage}{\linewidth}\raggedright{\scriptsize Notes. TV between the identified and naive one-step predictives of $\log P_{t+1}$, evaluated at the observed history. \texttt{Step frac} is the fraction of pooled steps with TV $\ge 0.99$; episode-level maxima can be below 1 even when some steps are at 1. The post phase is not reported in the current JSON. At the $\tau=2$ anchor the two predictives are maximally different (TV = 1); during the freeze TV is 0.65 ($\tau=2$) and 0.77 ($\tau=81$).}\end{minipage}\par

\par\textbf{Panel C. Policy boundary/action-flip diagnostic}

\par\textbf{Panel C.1 Flip probability}

\begin{adjustbox}{max width=\linewidth}
\begin{tabular}{@{}llllllllll@{}}
\toprule
Policy & $\tau$ (h) & $n$ & Flip mean & Flip SD & Flip p90 & Flip max & Mean $p_{\rm id}$ & Mean $p_{\rm naive}$ & Frac flip \textgreater{} 0.5 \\
\midrule
Wide 2\% & $\tau=2$ & 128 & 0.131 & 0.081 & 0.248 & 0.329 & 0.201 & 0.073 & 0.000 \\
Wide 2\% & $\tau=81$ & 128 & 0.456 & 0.239 & 0.737 & 0.814 & 0.737 & 0.361 & 0.531 \\
Narrow 1\% & $\tau=2$ & 128 & 0.214 & 0.092 & 0.334 & 0.433 & 0.460 & 0.247 & 0.000 \\
Narrow 1\% & $\tau=81$ & 128 & 0.318 & 0.254 & 0.700 & 0.793 & 0.865 & 0.644 & 0.328 \\
\bottomrule
\end{tabular}
\end{adjustbox}
\vspace{0.5em}
\par\textbf{Panel C.2 Boundary proximity (context)}

\begin{adjustbox}{max width=\linewidth}
\begin{tabular}{@{}llllllll@{}}
\toprule
Policy & $\tau$ (h) & $n$ & Mean margin & Frac rebalance & Frac $\vert m\vert\le0.0025$ & Frac $\vert m\vert\le0.005$ & Frac $\vert m\vert\le0.01$ \\
\midrule
Wide 2\% & $\tau=2$ & 128 & -0.0126 & 0.062 & 0.031 & 0.109 & 0.281 \\
Wide 2\% & $\tau=81$ & 128 & -0.0093 & 0.133 & 0.117 & 0.219 & 0.469 \\
Narrow 1\% & $\tau=2$ & 128 & -0.0034 & 0.234 & 0.219 & 0.453 & 0.969 \\
Narrow 1\% & $\tau=81$ & 128 & -0.0021 & 0.281 & 0.250 & 0.469 & 0.953 \\
\bottomrule
\end{tabular}
\end{adjustbox}
\vspace{0.5em}
\par\noindent\begin{minipage}{\linewidth}\raggedright{\scriptsize Notes. Flip = minimal-coupling probability that the rebalance decision at $t+1$ differs, $|p_{\rm id}-p_{\rm naive}|$, where $p_{\rm model} = P(|P_{t+1}/{\rm anchor}-1| > \theta)$ under that model's predictive; positions are from the observed trajectory (first-order). $p_{\rm id}$ and $p_{\rm naive}$ are the model-implied rebalance probabilities. Margin = $|P_t/{\rm anchor}-1| - \theta$ at decision steps; $m>0$ means the observed path rebalances. The boundary-proximity panel is context for the flip diagnostic, not a fourth term. The flip diagnostic is not a bound: it describes how close the policy is to its decision boundary under the two models.}\end{minipage}\par

\end{table*}

\begin{table*}[!htbp]
\centering
\scriptsize
\caption{Non-Markov Brownian bridge stress}
\label{tab:b2}
\vspace{0.5em}
\par\textbf{Panel A. Coverage and width}

\begin{adjustbox}{max width=\linewidth}
\begin{tabular}{@{}lllllllll@{}}
\toprule
Environment & Tier & open-only cov. (SE) & open width & cross-event cov. (SE) & cross width & width ratio & $\tau$-stratified cov. (SE) & $\tau$-stratified width ratio \\
\midrule
$\nu_{\rm mult}=0$ placebo & oracle & 0.897 (0.008) & 0.1618 & 0.897 (0.008) & 0.1637 & 1.01$\times$ & 0.901 (0.008) & 1.01$\times$ \\
$\nu_{\rm mult}=0$ placebo & estimated & 0.882 (0.008) & 0.1576 & 0.900 (0.008) & 0.1662 & 1.05$\times$ & 0.900 (0.008) & 1.05$\times$ \\
$\nu_{\rm mult}=0.6$ bridge & oracle & 0.818 (0.010) & 0.1618 & 0.900 (0.008) & 0.2047 & 1.27$\times$ & 0.902 (0.008) & 1.27$\times$ \\
$\nu_{\rm mult}=0.6$ bridge & estimated & 0.806 (0.010) & 0.1571 & 0.900 (0.008) & 0.2072 & 1.32$\times$ & 0.901 (0.008) & 1.31$\times$ \\
\bottomrule
\end{tabular}
\end{adjustbox}
\vspace{0.5em}
\par\noindent\begin{minipage}{\linewidth}\raggedright{\scriptsize Notes. 90\% prediction intervals, 5-fold out-of-fold (fold = seed mod 5), 128 seeds $\times$ 4 windows = 512 paths, 1,536 double-endpoint events ($\tau=2/65/81$, 512 each). Nominal 0.90 ($\alpha=0.1$). Coverage SEs are binomial over the 1,536 event predictions and ignore within-path and across-fold dependence, so they are lower bounds. \texttt{open-only} uses ordinary OPEN$\to$OPEN one-step transitions and the Markov $\tau$-step variance $v_{0e}$; \texttt{cross-event} uses $\operatorname{Var}(r_e)=v_{0e}+\theta\sigma_e^2\,\mathrm{gap}_e$ with $\theta$ fitted by Gaussian MLE; \texttt{$\tau$-stratified} uses a separate conformal quantile per freeze length and no functional form in $\tau$.}\end{minipage}\par

\par\textbf{Panel B. Bridge variance component}

\begin{adjustbox}{max width=\linewidth}
\begin{tabular}{@{}llllll@{}}
\toprule
Environment & Tier & $\hat\theta$ & 95\% CI & $\hat\nu/\sigma$ & 95\% CI \\
\midrule
$\nu_{\rm mult}=0$ placebo & oracle & 0.0282 & [0.000, 0.097] & 0.168 & [0.000, 0.311] \\
$\nu_{\rm mult}=0$ placebo & estimated & 0.0703 & [0.008, 0.138] & 0.265 & [0.092, 0.372] \\
$\nu_{\rm mult}=0.6$ bridge & oracle & 0.4912 & [0.397, 0.595] & 0.701 & [0.630, 0.771] \\
$\nu_{\rm mult}=0.6$ bridge & estimated & 0.5360 & [0.443, 0.639] & 0.732 & [0.665, 0.799] \\
\bottomrule
\end{tabular}
\end{adjustbox}
\vspace{0.5em}
\par\noindent\begin{minipage}{\linewidth}\raggedright{\scriptsize Notes. $\theta=\nu_{\rm mult}^2\,\mathbb E[e^{2\zeta}]$; the design value is $\theta_{\rm true}=0.457$ from the realized event-level draws. The $\nu=0$ placebo has $\theta=0$ by construction. The oracle-tier placebo interval contains zero; the estimated-tier placebo has a small positive bias from fitting $T$ and $\sigma$ on open transitions, as in the $\gamma_{\mathrm{st}}=0$ placebo of Appendix Table B10.}\end{minipage}\par

\par\textbf{Panel C. Coverage by freeze length}

\begin{adjustbox}{max width=\linewidth}
\begin{tabular}{@{}llllllllll@{}}
\toprule
Environment & Tier & $\tau$ & $n$ & open cov. & open width & cross cov. & cross width & $\tau$-strat cov. & $\tau$-strat width \\
\midrule
$\nu_{\rm mult}=0$ placebo & oracle & 2 & 512 & 0.891 & 0.0488 & 0.893 & 0.0492 & 0.900 & 0.0497 \\
$\nu_{\rm mult}=0$ placebo & oracle & 65 & 512 & 0.891 & 0.2089 & 0.891 & 0.2113 & 0.902 & 0.2188 \\
$\nu_{\rm mult}=0$ placebo & oracle & 81 & 512 & 0.910 & 0.2277 & 0.908 & 0.2305 & 0.900 & 0.2239 \\
$\nu_{\rm mult}=0$ placebo & estimated & 2 & 512 & 0.887 & 0.0479 & 0.902 & 0.0502 & 0.898 & 0.0499 \\
$\nu_{\rm mult}=0$ placebo & estimated & 65 & 512 & 0.865 & 0.2034 & 0.893 & 0.2145 & 0.902 & 0.2213 \\
$\nu_{\rm mult}=0$ placebo & estimated & 81 & 512 & 0.895 & 0.2213 & 0.906 & 0.2339 & 0.900 & 0.2273 \\
$\nu_{\rm mult}=0.6$ bridge & oracle & 2 & 512 & 0.820 & 0.0488 & 0.895 & 0.0597 & 0.904 & 0.0612 \\
$\nu_{\rm mult}=0.6$ bridge & oracle & 65 & 512 & 0.797 & 0.2089 & 0.889 & 0.2641 & 0.900 & 0.2800 \\
$\nu_{\rm mult}=0.6$ bridge & oracle & 81 & 512 & 0.838 & 0.2277 & 0.916 & 0.2903 & 0.900 & 0.2744 \\
$\nu_{\rm mult}=0.6$ bridge & estimated & 2 & 512 & 0.812 & 0.0479 & 0.902 & 0.0607 & 0.900 & 0.0610 \\
$\nu_{\rm mult}=0.6$ bridge & estimated & 65 & 512 & 0.791 & 0.2029 & 0.879 & 0.2672 & 0.902 & 0.2879 \\
$\nu_{\rm mult}=0.6$ bridge & estimated & 81 & 512 & 0.814 & 0.2206 & 0.920 & 0.2935 & 0.900 & 0.2708 \\
\bottomrule
\end{tabular}
\end{adjustbox}
\vspace{0.5em}
\par\noindent\begin{minipage}{\linewidth}\raggedright{\scriptsize Notes. The open-only deficit is present at every $\tau$; it is not a single long-freeze cell. The $\tau$-stratified calibrator recovers within-$\tau$ marginal coverage without assuming the bridge variance shape.}\end{minipage}\par

\par\textbf{Panel D. Same endpoints, different bridge volatility}

\begin{adjustbox}{max width=\linewidth}
\begin{tabular}{@{}ll@{}}
\toprule
Quantity & Value \\
\midrule
Same-endpoint events & 64 \\
Max abs pool-price difference & 0.0 \\
Observed paths bit-identical & True \\
Mean latent QV at $\nu_e$ & 1.6023e-03 \\
Mean latent QV at $2\nu_e$ & 6.1677e-03 \\
Empirical QV ratio & 3.85 \\
Theoretical QV ratio & 4.00 \\
\bottomrule
\end{tabular}
\end{adjustbox}
\vspace{0.5em}
\par\noindent\begin{minipage}{\linewidth}\raggedright{\scriptsize Notes. Each selected event is re-simulated twice with the same realized endpoints, the same standard Brownian bridge shape, and volatilities $\nu_e$ and $2\nu_e$. Freeze interiors are masked by the observation map, so the observed pool-price paths are bit-identical; only the latent quadratic variation changes. The empirical QV ratio is close to the theoretical $2^2=4$. This is the empirical counterpart of the Theorem 4 construction, not a coverage experiment: with endpoints pinned, no observable-data estimator can distinguish the two bridge laws.}\end{minipage}\par

\par\textbf{Panel E. Coverage against the realized event-level volatility}

\begin{adjustbox}{max width=\linewidth}
\begin{tabular}{@{}llllllll@{}}
\toprule
$\nu_e$ quartile & $n$ & $\nu$ min & $\nu$ max & open cov. & cross cov. & open width & cross width \\
\midrule
0 & 384 & 0.0010 & 0.0034 & 0.859 & 0.919 & 0.1210 & 0.1534 \\
1 & 384 & 0.0034 & 0.0048 & 0.844 & 0.917 & 0.1451 & 0.1839 \\
2 & 384 & 0.0048 & 0.0064 & 0.797 & 0.904 & 0.1654 & 0.2093 \\
3 & 384 & 0.0064 & 0.0293 & 0.773 & 0.859 & 0.2157 & 0.2724 \\
\bottomrule
\end{tabular}
\end{adjustbox}
\vspace{0.5em}
\par\noindent\begin{minipage}{\linewidth}\raggedright{\scriptsize Notes. The groups use the simulation-only $\nu_e$, which is not observed by either method; they are a dose-response diagnostic, not a deployable stratification. Open-only coverage falls monotonically across all four quartiles; cross-event calibration recovers marginal coverage but not coverage conditional on the latent $\nu_e$ (the highest quartile remains below nominal). The latter is expected: $\nu_e$ is not in the observed information set.}\end{minipage}\par

\par\noindent\begin{minipage}{\linewidth}\raggedright{\scriptsize Bridge law Brownian bridge pinned at the latent pre/post anchors. The post endpoint is drawn from $N(m_+, v_{\rm ar}+\nu_e^2\,\mathrm{gap})$, with $\nu_e=\sigma\nu_{\rm mult}e^{\zeta_e}$, $\zeta_e\sim N(0,0.35^2)$ clipped to $[-1.0,1.0]$, and $\nu_{\rm mult}=0.60$. The interior is unobserved, so the bridge volatility's observable footprint is the endpoint law; the same-endpoint panel isolates its path-functional consequence.}\end{minipage}\par

\end{table*}

\begin{table*}[!htbp]
\centering
\scriptsize
\caption{Semi-synthetic calibration}
\label{tab:b3}
\par\noindent\begin{minipage}{\linewidth}\raggedright{\scriptsize Real hourly bars (Yahoo Finance, 2025-10-22 to 2026-10-21, pre/post included) calibrate the return dynamics; freeze durations are measured Chainlink stale intervals. The full observation map is not independently calibrated by these inputs.}\end{minipage}\par

\par\textbf{Panel A. Return dynamics (regular session)}

\begin{adjustbox}{max width=\linewidth}
\begin{tabular}{@{}lllllll@{}}
\toprule
Symbol & $n$ bars & $\sigma_h$ ($10^{-3}$) & $\sigma_h^{24/7}$ ($10^{-3}$) & $\hat\kappa_h$ ($10^{-3}$) & jump prob. ($\lvert r\rvert>4\sigma_h$) & jump sd / $\sigma_h$ \\
\midrule
SPY & 4147 & 2.21 & 1.47 & 0.90 & 0.0060 & 4.99 \\
QQQ & 4128 & 3.02 & 2.16 & 0.71 & 0.0060 & 4.18 \\
TSLA & 4128 & 6.51 & 4.86 & 2.01 & 0.0033 & 5.73 \\
NVDA & 4128 & 5.77 & 4.23 & 3.61 & 0.0047 & 5.08 \\
\bottomrule
\end{tabular}
\end{adjustbox}
\vspace{0.5em}
\par\textbf{Panel B. Freeze duration and gap benchmarks}

\begin{adjustbox}{max width=\linewidth}
\begin{tabular}{@{}lllllllll@{}}
\toprule
Symbol & \makecell{weekend $\tau$\\median (h)} & \makecell{weekend $\tau$\\p90 (h)} & \makecell{holiday $\tau$\\median (h)} & \makecell{holiday $\tau$\\p90 (h)} & \makecell{weekend gap /\\$\sigma_h\sqrt{\tau}$} & \makecell{holiday gap /\\$\sigma_h\sqrt{\tau}$} & \makecell{weekend gap /\\$g(\gamma_{\mathrm{st}}=4)$} & \makecell{holiday gap /\\$g(\gamma_{\mathrm{st}}=4)$} \\
\midrule
SPY & 56.1 & 60.7 & 78.7 & 80.2 & 0.35 & 0.40 & 0.77 & 0.95 \\
QQQ & 55.4 & 58.3 & 76.3 & 87.5 & 0.40 & 0.38 & 0.86 & 0.90 \\
TSLA & 52.1 & 54.7 & 76.1 & 78.0 & 0.26 & 0.27 & 0.54 & 0.65 \\
NVDA & 52.1 & 55.7 & 76.0 & 77.6 & 0.28 & 0.29 & 0.59 & 0.70 \\
\bottomrule
\end{tabular}
\end{adjustbox}
\vspace{0.5em}
\par\textbf{Panel C. Observation-map parameters (not independently calibrated here)}

\begin{adjustbox}{max width=\linewidth}
\begin{tabular}{@{}llllll@{}}
\toprule
Source & $\sigma_o$ & $\phi_{ext}$ & $\sigma_{ext}$ & $\phi_{stale}$ & $\sigma_{stale}$ \\
\midrule
G0 true-$S$ diagnostic & 0.00050 & 0.805 & 0.00202 & 0.9784 & 0.00080 \\
observable-only estimate & 0.00041 & 0.774 & 0.00202 & 0.9802 & 0.00080 \\
\bottomrule
\end{tabular}
\end{adjustbox}
\vspace{0.5em}
\par\noindent\begin{minipage}{\linewidth}\raggedright{\scriptsize Notes. $\sigma_h$ is the standard deviation of hourly log returns within the regular session (09:30--16:00 ET, gaps $\le 90$ min); bars are stamped at the start of their interval, so the 16:00 bar is the first after-hours bar and is excluded. $\sigma_h^{24/7}=\sqrt{\sum r^2/\sum\Delta t}$ over all consecutive bars including overnight/weekend/holiday gaps, i.e. a per-elapsed-hour realised variance rate; $\hat\kappa_h=1-\hat\phi$ from an AR(1) on consecutive regular-session bars. Jump probability and conditional sd use the $4\sigma_h$ tail (5--9 events per symbol, so the conditional sd is noisy); the conditional sd is descriptive and not a jump-diffusion MLE. Overnight/weekend/holiday sd are log-close changes across duration buckets of 6--40 h, 40--70 h, and $>70$ h; the 40--70 h bucket can include midweek holiday closures, since it is not a calendar label. Freeze durations are measured Chainlink stale intervals for the matching 24/5 feed: weekend rows with gap $\ge 40$ h (28 gaps per feed), and the longest gap per labelled holiday (5 holidays per feed) before taking the median/p90 across holidays (short partial gaps inherit the calendar label and would otherwise bias the median down). The measured weekend stale interval (52--56 h) is shorter than the Friday-close-to-Monday-open market gap (about 65 h) because the 24/5 feed resumes on Sunday evening. In Panel B, gap/$\sigma_h\sqrt{\tau}$ is the continuous-diffusion benchmark, and gap/$g(\gamma_{\mathrm{st}}=4)$ compares the observed total gap sd with the jump-only scale of the paper's stress variant ($g=\gamma\sigma_h\sqrt{1-e^{-\tau/40}}$, $\gamma_{\mathrm{st}}=4$, $\tau_0=40$). The gap sd is measured across the market closure (about 65 h on weekends), while $\tau$ is the oracle stale interval (52--56 h); using the longer market duration in the denominator would lower the ratios, so the reported values are conservative. Values below 1 mean the observed gap is smaller than the benchmark. Because the observed gap also contains any freeze-period latent movement, the jump component itself is even smaller; the stress design is therefore conservative relative to these feeds. The observation map $h$ is not independently estimated by these calibration inputs: public bars provide return benchmarks and feed logs provide stale schedules, rather than a paired sample for fitting the full map; both synthetic observation-map estimates are reported, and the paper's coverage experiments use the observable-only one as their main specification. Calibration summary. The simulator's DGP uses $\sigma_h=8.0\times10^{-3}$, $\kappa_h=3.0\times10^{-3}$, jump probability 0.005 with jump sd $3\sigma_h$, weekend freeze 65 h, and holiday freeze 81 h. The four feeds span $\sigma_h=2.21$--$6.51\times10^{-3}$, $\sigma_h^{24/7}=1.47$--$4.86\times10^{-3}$ (annualised 14--45\%), $\kappa_h=0.71$--$3.61\times10^{-3}$, jump probability 0.0033--0.0060 (bracketing the simulator's 0.005) with conditional sd 4.2$\sigma_h$--5.7$\sigma_h$, weekend freeze 52.1--56.1 h, and holiday freeze 76.0--78.7 h. The simulator's $\sigma_h=8.0\times10^{-3}$ is above every session estimate and 1.6--5.4$\times$ the 24/7-equivalent rate (its annualised latent vol is about 75\%); the environment is therefore a high-volatility stress setting, not a broad-market calibration. Its mean reversion sits inside the measured range ($0.71$--$3.61\times10^{-3}$), close to NVDA's upper end; its weekend freeze is longer than the measured 52--56 h; its 81 h holiday freeze is close to the measured 76--79 h; and its 3$\sigma_h$ jump sd is thinner-tailed than the measured 4.2$\sigma_h$--5.7$\sigma_h$ conditional sd.}\end{minipage}\par

\end{table*}

\begin{table*}[!htbp]
\centering
\scriptsize
\caption{Misspecification suite}
\label{tab:b4}
\vspace{0.5em}
\par\textbf{Panel A. Wide policy: coverage under one perturbation magnitude per class}

\begin{adjustbox}{max width=\linewidth}
\begin{tabular}{@{}lllllllll@{}}
\toprule
Variant & Model & bias $G$ & $G$ within & $G$ LORO & $G_r$ within & $G_r$ LORO & plug-in $G$ & $G$ LORO width \\
\midrule
Correctly specified & identified & -0.0004 & 0.932 & 0.924 & 0.909 & 0.871 & 0.852 & 0.0463 \\
Correctly specified & naive & +0.0087 & 0.912 & 0.904 & 0.898 & 0.875 & 0.662 & 0.0448 \\
AR(2) observation noise & identified & -0.0004 & 0.916 & 0.912 & 0.903 & 0.875 & 0.846 & 0.0458 \\
AR(2) observation noise & naive & +0.0089 & 0.912 & 0.898 & 0.894 & 0.873 & 0.652 & 0.0438 \\
Student-$t(3)$ jumps & identified & +0.0002 & 0.920 & 0.910 & 0.895 & 0.875 & 0.848 & 0.0436 \\
Student-$t(3)$ jumps & naive & +0.0087 & 0.906 & 0.898 & 0.900 & 0.874 & 0.621 & 0.0441 \\
Regime-label noise & identified & -0.0019 & 0.904 & 0.900 & 0.900 & 0.870 & 0.861 & 0.0440 \\
Regime-label noise & naive & +0.0086 & 0.912 & 0.902 & 0.900 & 0.873 & 0.643 & 0.0456 \\
Hidden stale state & identified & +0.0004 & 0.928 & 0.924 & 0.891 & 0.863 & 0.846 & 0.0460 \\
Hidden stale state & naive & +0.0076 & 0.893 & 0.891 & 0.902 & 0.872 & 0.693 & 0.0410 \\
\bottomrule
\end{tabular}
\end{adjustbox}
\vspace{0.5em}
\par\textbf{Panel B. Narrow policy}

\begin{adjustbox}{max width=\linewidth}
\begin{tabular}{@{}lllllllll@{}}
\toprule
Variant & Model & bias $G$ & $G$ within & $G$ LORO & $G_r$ within & $G_r$ LORO & plug-in $G$ & $G$ LORO width \\
\midrule
Correctly specified & identified & -0.0017 & 0.914 & 0.900 & 0.902 & 0.866 & 0.826 & 0.1061 \\
Correctly specified & naive & +0.0218 & 0.900 & 0.887 & 0.898 & 0.881 & 0.668 & 0.1196 \\
AR(2) observation noise & identified & -0.0028 & 0.920 & 0.914 & 0.903 & 0.863 & 0.838 & 0.1113 \\
AR(2) observation noise & naive & +0.0216 & 0.908 & 0.896 & 0.900 & 0.874 & 0.666 & 0.1254 \\
Student-$t(3)$ jumps & identified & -0.0004 & 0.922 & 0.914 & 0.900 & 0.874 & 0.826 & 0.1094 \\
Student-$t(3)$ jumps & naive & +0.0214 & 0.900 & 0.885 & 0.902 & 0.873 & 0.633 & 0.1210 \\
Regime-label noise & identified & -0.0053 & 0.906 & 0.914 & 0.902 & 0.868 & 0.832 & 0.1109 \\
Regime-label noise & naive & +0.0219 & 0.910 & 0.889 & 0.901 & 0.884 & 0.670 & 0.1210 \\
Hidden stale state & identified & +0.0003 & 0.912 & 0.912 & 0.898 & 0.861 & 0.822 & 0.1131 \\
Hidden stale state & naive & +0.0180 & 0.912 & 0.895 & 0.896 & 0.878 & 0.723 & 0.1189 \\
\bottomrule
\end{tabular}
\end{adjustbox}
\vspace{0.5em}
\par\textbf{Panel C. Paired identified $-$ naive coverage}

\begin{adjustbox}{max width=\linewidth}
\begin{tabular}{@{}lllll@{}}
\toprule
Variant & Policy & $G$ studentized & $G_r$ studentized & plug-in $G$ \\
\midrule
Correctly specified & Wide & +0.020 [-0.014, +0.053] & +0.012 [-0.003, +0.027] & +0.189 [+0.141, +0.240] \\
Correctly specified & Narrow & -0.037 [-0.070, -0.004] & +0.017 [+0.003, +0.031] & +0.158 [+0.104, +0.209] \\
AR(2) observation noise & Wide & +0.041 [+0.008, +0.074] & +0.017 [+0.003, +0.031] & +0.193 [+0.141, +0.246] \\
AR(2) observation noise & Narrow & +0.010 [-0.023, +0.043] & +0.008 [-0.007, +0.022] & +0.172 [+0.121, +0.223] \\
Student-$t(3)$ jumps & Wide & -0.037 [-0.068, -0.006] & +0.015 [+0.001, +0.029] & +0.227 [+0.172, +0.281] \\
Student-$t(3)$ jumps & Narrow & +0.008 [-0.029, +0.045] & +0.018 [+0.003, +0.034] & +0.193 [+0.137, +0.246] \\
Regime-label noise & Wide & +0.020 [-0.014, +0.055] & +0.002 [-0.012, +0.016] & +0.219 [+0.170, +0.271] \\
Regime-label noise & Narrow & -0.010 [-0.043, +0.021] & -0.005 [-0.020, +0.010] & +0.162 [+0.111, +0.211] \\
Hidden stale state & Wide & +0.033 [-0.006, +0.068] & -0.000 [-0.015, +0.014] & +0.152 [+0.100, +0.203] \\
Hidden stale state & Narrow & +0.029 [-0.004, +0.062] & +0.016 [-0.000, +0.032] & +0.100 [+0.051, +0.148] \\
\bottomrule
\end{tabular}
\end{adjustbox}
\vspace{0.5em}
\par\noindent\begin{minipage}{\linewidth}\raggedright{\scriptsize Notes. 128 calibration and 128 evaluation seeds per window, four windows (512 calibration and 512 evaluation paths), 20 model draws per path; nominal 90\% ($\alpha=0.1$). $G$ = full window, $G_r$ = regime segment; within = same-population calibration, LORO = leave-one-window-out ($G$) or leave-one-regime-out ($G_r$). Bias is the mean model-mean minus truth over evaluation paths, in units of the initial position value $V_0=10^5$ USD ($0.001=\$100$). Panel C reports identified minus naive with a 95\% path bootstrap interval (2,000 resamples); $G$ and $G_r$ use the studentized score. The model side is fixed at the G0 $\kappa,\mu$, jump probability, and $h$ for every variant; only the data-generating process changes. \texttt{AR(2) observation noise} replaces the EXTENDED AR(1) with $(\phi_1,\phi_2)=(0.70,0.20)$ and the stale AR(1) with $(0.90,0.08)$, matching the stationary sd. \texttt{Student-$t(3)$ jumps} keeps the 0.005 jump probability and scales the jump to sd $3\sigma$. \texttt{Regime-label noise} flips 15\% of OPEN and EXTENDED labels; the latent path and $P$ are unchanged. \texttt{Hidden stale state} adds a simulation-only two-state chain during dense hours (enter 0.01, exit 0.125 per hour; stationary occupancy 7.4\%) that switches the observation law to the stale map; the model and policy still see the calendar.}\end{minipage}\par

\par\textbf{Panel D. Design}

\begin{adjustbox}{max width=\linewidth}
\begin{tabular}{@{}ll@{}}
\toprule
Ingredient & Setting \\
\midrule
AR(2) EXTENDED & $\phi=(0.70, 0.20)$, stationary sd 0.0033 \\
AR(2) stale & $\phi=(0.90, 0.08)$, stationary sd 0.0040 \\
Student-$t(3)$ jumps & jump probability 0.005; sd $3\sigma$ \\
Regime-label noise & flip probability 0.15 on OPEN/EXTENDED \\
Hidden stale state & enter 0.010, exit 0.125, stationary occupancy 7.4\%, burn-in 24 h \\
\bottomrule
\end{tabular}
\end{adjustbox}
\vspace{0.5em}
\end{table*}

\begin{table}[!ht]
\centering
\scriptsize
\caption{Row-level and trajectory-level coverage}
\label{tab:3}
\vspace{0.5em}
\par\textbf{Panel A. L0: held-out-regime row-level coverage (90\% split conformal)}

\begin{adjustbox}{max width=\linewidth}
\begin{tabular}{@{}llllllll@{}}
\toprule
Regime & $n_{\rm test}$ & Naive global 90\% cov. [95\% CI] & width & Mondrian cov. & width ratio & Similarity cov. & width ratio \\
\midrule
OPEN & 46,848 & 0.9239 [0.9215, 0.9263] & 0.00614 & 0.9349 & 1.12$\times$ & 0.9033 & 0.83$\times$ \\
EXTENDED & 97,536 & 0.9301 [0.9285, 0.9317] & 0.00644 & 0.9378 & 1.08$\times$ & 0.9055 & 0.79$\times$ \\
WEEKEND & 92,928 & 0.8626 [0.8604, 0.8648] & 0.00517 & 0.9020 & 1.31$\times$ & 0.9020 & 1.31$\times$ \\
HOLIDAY & 18,432 & 0.8990 [0.8947, 0.9034] & 0.00597 & 0.9156 & 1.15$\times$ & 0.9156 & 1.15$\times$ \\
HALT & 1,536 & 0.8906 [0.8750, 0.9062] & 0.00598 & 0.9017 & 1.15$\times$ & 0.9017 & 1.15$\times$ \\
\bottomrule
\end{tabular}
\end{adjustbox}
\vspace{0.5em}
\par\noindent\begin{minipage}{\linewidth}\raggedright{\scriptsize Notes. Nominal coverage is 0.90 ($\alpha=0.1$); the 95\% CI is a binomial normal approximation on $n_{\rm test}$ rows and ignores within-trajectory correlation, so it is a lower bound on the true uncertainty. Widths are mean interval widths. Naive global = one pooled calibration radius (baseline). Mondrian = one radius per calibration regime, unseen regime falls back to the largest calibration radius. Similarity = Mondrian with the two-class fallback \{OPEN, EXTENDED\} / \{WEEKEND, HOLIDAY, HALT\}. Baseline aggregate \texttt{coverage\_gap} = 0.0136; similarity fallback = 0.0000; strategy-conditional = 0.00076. Strategy-weighted gap: naive 0.0429, similarity 0.0392, strategy-conditional 0.0024. WEEKEND coverage by window (W1--W4): W1 0.8550, W2 0.8782, W3 0.8706, W4 0.8467. HALT calibration radius: 384 rows per fold, row-bootstrap SE 0.000296--0.000315 (8.8--9.3\% of the radius).}\end{minipage}\par

\par\textbf{Panel B. L1: trajectory-level 90\% conformal coverage and width}

\begin{adjustbox}{max width=\linewidth}
\begin{tabular}{@{}lllllllll@{}}
\toprule
Policy & Model & Score & $G$ within & $G$ LORO & $G_r$ within & $G_r$ LORO & width $G$ & width $G_r$ \\
\midrule
Wide & identified & raw & 0.9297 & 0.9238 & 0.9062 & 0.8719 & 0.04752 & 0.01585 \\
Wide & identified & studentized & 0.9180 & 0.9121 & 0.9012 & 0.8965 & 0.04158 & 0.01373 \\
Wide & identified & cqr & 0.9102 & 0.9082 & 0.9035 & 0.9004 & 0.03867 & 0.01229 \\
Wide & naive & raw & 0.9199 & 0.9102 & 0.9000 & 0.8738 & 0.04452 & 0.01823 \\
Wide & naive & studentized & 0.9121 & 0.9102 & 0.8949 & 0.8812 & 0.04890 & 0.01411 \\
Wide & naive & cqr & 0.9141 & 0.9004 & 0.8930 & 0.8836 & 0.04138 & 0.01342 \\
Wide & identified & plug-in normal & 0.8438 & -- & 0.8887 & -- & 0.03586 & 0.01283 \\
Wide & naive & plug-in normal & 0.6660 & -- & 0.7336 & -- & 0.02669 & 0.00895 \\
Narrow & identified & raw & 0.9199 & 0.9121 & 0.9059 & 0.8668 & 0.10866 & 0.04785 \\
Narrow & identified & studentized & 0.8984 & 0.8926 & 0.9066 & 0.8875 & 0.10100 & 0.03806 \\
Narrow & identified & cqr & 0.8945 & 0.8906 & 0.9082 & 0.9008 & 0.09946 & 0.03647 \\
Narrow & naive & raw & 0.9082 & 0.8906 & 0.8980 & 0.8801 & 0.12040 & 0.05784 \\
Narrow & naive & studentized & 0.9277 & 0.9258 & 0.8918 & 0.8801 & 0.13524 & 0.04994 \\
Narrow & naive & cqr & 0.9141 & 0.9062 & 0.8965 & 0.8699 & 0.11310 & 0.04547 \\
Narrow & identified & plug-in normal & 0.8359 & -- & 0.8820 & -- & 0.08656 & 0.03498 \\
Narrow & naive & plug-in normal & 0.6699 & -- & 0.7113 & -- & 0.07261 & 0.02774 \\
\bottomrule
\end{tabular}
\end{adjustbox}
\vspace{0.5em}
\par\noindent\begin{minipage}{\linewidth}\raggedright{\scriptsize Notes. $G$ is the full-window cumulative reward and $G_r$ the regime-segment reward; within = calibrate and test on the same regime/window population; LORO = leave-one-regime-out (for $G_r$) or leave-one-window-out (for $G$). Widths are mean 90\% interval widths: for conformal rows they are the LORO widths, whereas for plug-in normal rows they are the full-window ($G$) and within-regime ($G_r$) widths, since the plug-in has no LORO variant. Plug-in normal = $\pm z_{0.95}\,\hat{\rm sd}$, not a conformal interval. Identified uses the regime-dependent observation model with G0 true-$S$ diagnostic parameters; naive treats $P$ as $S$ with per-path $\kappa_P/\mu_P/\sigma_{\rm all}$.}\end{minipage}\par

\par\textbf{Panel C. L1: path-level paired bootstrap, identified $-$ naive}

\begin{adjustbox}{max width=\linewidth}
\begin{tabular}{@{}lllll@{}}
\toprule
Policy & Level & Score & identified $-$ naive & 95\% CI \\
\midrule
Wide & full-window & studentized & +0.0020 & [-0.0312, +0.0352] \\
Wide & full-window & cqr & +0.0078 & [-0.0254, +0.0430] \\
Wide & full-window & plugin\_normal & +0.1777 & [+0.1270, +0.2285] \\
Wide & regime & studentized & +0.0152 & [+0.0012, +0.0289] \\
Wide & regime & cqr & +0.0168 & [+0.0031, +0.0312] \\
Wide & regime & plugin\_normal & +0.1551 & [+0.1379, +0.1730] \\
Narrow & full-window & studentized & -0.0332 & [-0.0684, +0.0020] \\
Narrow & full-window & cqr & -0.0156 & [-0.0508, +0.0215] \\
Narrow & full-window & plugin\_normal & +0.1660 & [+0.1152, +0.2168] \\
Narrow & regime & studentized & +0.0074 & [-0.0074, +0.0223] \\
Narrow & regime & cqr & +0.0309 & [+0.0164, +0.0461] \\
Narrow & regime & plugin\_normal & +0.1707 & [+0.1520, +0.1891] \\
\bottomrule
\end{tabular}
\end{adjustbox}
\vspace{0.5em}
\par\noindent\begin{minipage}{\linewidth}\raggedright{\scriptsize Notes. 2,000 path-level resamples, conditioning on the fitted calibration quantiles; the interval quantifies test-path sampling only.}\end{minipage}\par

\end{table}

% Table 9 --- mirrored to paper/main-draft.md "Table 9".
% Rows 2--6 are title-level descriptions backed by P / metadata-only
% citations; they carry no method detail, comparative conclusion, or
% exclusivity claim.  Rows 1 and 7 are supported at the cited theorem
% or section level (paper/citation-map-m4.md).
\begin{table*}[!ht]
\centering
\footnotesize
\setlength{\tabcolsep}{3pt}
\caption{Positioning against related evaluation objects.}
\label{tab:9}
\begin{tabular}{@{}>{\raggedright\arraybackslash}p{0.030\textwidth}
  >{\raggedright\arraybackslash}p{0.175\textwidth}
  >{\raggedright\arraybackslash}p{0.290\textwidth}
  >{\raggedright\arraybackslash}p{0.455\textwidth}@{}}
\toprule
\textbf{\#} & \textbf{Line of work} & \textbf{Topic studied} &
\textbf{This paper} \\
\midrule
1 & Mondrian / group conformal \cite{vovk2005,lei2018} &
exchangeable units; group label known before prediction; finite-sample
within-group coverage &
the branch label is post-decision; for branch-blind, nondegenerate score rules
Theorem~3 shows a pre-decision Mondrian partition cannot recover branch coverage;
Proposition~2 prices calibration-side labels \\
\addlinespace
2 & Weighted conformal \cite{tibshirani2019} &
conformal prediction under covariate shift &
the shift is in the observation kernel, not the covariate law; the target
includes branch coverage \\
\addlinespace
3 & Adaptive conformal inference \cite{gibbs2021} &
conformal inference under distribution shift &
offline OPE with passive logs and no online feedback; the target includes
branch coverage \\
\addlinespace
4 & Conformal policy evaluation \cite{gan2025} &
conformal inference for policy evaluation &
the observation kernel changes with the calendar regime; the target includes
branch-conditional coverage; we add an impossibility map \\
\addlinespace
5 & POMDP OPE \cite{zhang2024,zhang2025,luo2026} &
off-policy evaluation in partially observed systems &
the regime index changes the kernel; logging is passive; identification under
I1$'$/I2/I3, with I5/I5$'$ only for a non-Markov bridge; the output is a
coverage map \\
\addlinespace
6 & Selective labels \cite{lakkaraju2017} &
evaluation under selective labels &
the unavailable object is the branch label at the pre-decision
$\sigma$-field; the target is coverage; Theorem~3 gives an impossibility \\
\addlinespace
7 & Limits of conditional coverage \cite{barber2021} &
conditional coverage; marginal-to-conditional upgrade impossible without
distributional assumptions &
the obstruction is the timing of the branch label; branch-wise calibration is
possible with calibration-side labels, at a width price (Proposition~2) \\
\addlinespace
8 & \textbf{This paper} &
calendar-indexed observation kernels $h_r$; latent state $S$; passive logs &
Theorem~1 (identification), Theorem~2 (drift-aware marginal coverage),
Theorem~3 (branch impossibility), Proposition~2 (width price), Theorem~4
(within-class non-identification), L0--L3 ladder \\
\bottomrule
\end{tabular}
\end{table*}

\begin{table}[!ht]
\centering
\scriptsize
\caption{Second synthetic sensor-POMDP testbed}
\label{tab:13}
\vspace{0.5em}
\par\textbf{Panel A. Return prediction error and interval coverage}

\begin{adjustbox}{max width=\linewidth}
\begin{tabular}{@{}lllll@{}}
\toprule
Model & Mean absolute return error (G) & Plug-in coverage & Raw full-window coverage & Raw LORO outage coverage \\
\midrule
identified transition & 4.805 & 0.906 & 0.883 & 0.758 \\
fitted transition & 5.850 & 0.727 & 0.914 & 0.812 \\
\bottomrule
\end{tabular}
\end{adjustbox}
\vspace{0.5em}
\par\noindent\begin{minipage}{\linewidth}\raggedright{\scriptsize Notes. Mean absolute realized-return prediction error is the mean over evaluation paths of the absolute difference between the 64-rollout predictive mean and the realized target return $G$; it is not a value-bias estimate. These are descriptive point estimates without paired uncertainty intervals. Calibration and evaluation use 128 independent paths each. Truth and both predictive models use the same target-action rule and identified-belief filter, with action realizations following their simulated observation histories. Both use the known control coefficient $B=0.35$. The identified-transition predictor also knows $A$, $\sigma$, and the observation model; the fitted-transition comparator replaces only the autonomous rollout AR(1) coefficient and residual scale, estimated from burn-in data. The JSON retains \texttt{naive} as its internal key. This is not an observation-filter comparison or evidence of uniform model superiority.}\end{minipage}\par

\par\textbf{Panel B. Raw-score coverage by observation regime}

\begin{adjustbox}{max width=\linewidth}
\begin{tabular}{@{}lllll@{}}
\toprule
Regime & Known-trans. within & Known-trans. LORO & Fitted-trans. within & Fitted-trans. LORO \\
\midrule
normal & 0.938 & 0.953 & 0.906 & 0.992 \\
degraded & 0.859 & 0.805 & 0.898 & 0.836 \\
outage & 0.859 & 0.758 & 0.906 & 0.812 \\
\bottomrule
\end{tabular}
\end{adjustbox}
\vspace{0.5em}
\par\noindent\begin{minipage}{\linewidth}\raggedright{\scriptsize Notes. LORO calibrates on the other two regimes. The outage boundary is retained rather than pooled away.}\end{minipage}\par

\par\textbf{Panel C. Post-decision intervention-branch coverage (identified, raw)}

\begin{adjustbox}{max width=\linewidth}
\begin{tabular}{@{}lllll@{}}
\toprule
Rule & Overall coverage & Branch 0 coverage & Branch 1 coverage & Mean width (return units) \\
\midrule
marginal & 0.883 & 0.966 & 0.692 & 14.851 \\
predecision & 0.891 & 0.978 & 0.692 & 16.288 \\
branch\_union & 0.953 & 0.978 & 0.897 & 28.057 \\
\bottomrule
\end{tabular}
\end{adjustbox}
\vspace{0.5em}
\par\noindent\begin{minipage}{\linewidth}\raggedright{\scriptsize Notes. Calibration has 43 intervention events among 128 paths; evaluation has 39 among 128 (the other 89 are branch 0). Marginal intervention-branch SE is 0.074; branch-union SE is 0.049. The pre-decision rule uses three groups from the sign and magnitude of the last burn-in observation. Branch-union uses the larger calibration-side branch threshold, so the test-time branch label is not needed to form the interval. Mean width is the interval width in target-return (reward) units, not a coverage probability; the 1.89$\times$ figure in the main text is the branch-union/marginal width ratio.}\end{minipage}\par

\end{table}

\begin{table}[!ht]
\centering
\scriptsize
\caption{Answers to Q1--Q3}
\label{tab:8}
\begin{tabular}{@{}>{\raggedright\arraybackslash}p{0.21\linewidth}>{\raggedright\arraybackslash}p{0.52\linewidth}>{\raggedright\arraybackslash}p{0.21\linewidth}@{}}
\toprule
Question & Answer & Evidence \\
\midrule
Q1. What is identifiable from observable data? & Under I1$'$/I2/I3 (with A1, A2, A4--A7, T1, T4, and A3$'$ where applicable), the short-range transition, the stale/extended observation scales, and the value functional are identified. The Markov bridge needs no I5/I5$'$; the non-Markov bridge additionally needs I5/T2/T3 for the endpoint law and I5$'$ for the path law, so its value conclusion carries those extra assumptions. $h_{\rm open,sd}/h_{\rm ext,\phi}$ are only weakly identified. & \S3, Theorem 1; \S7.5, Table~\ref{tab:5} \\
Q2. What is recoverable by calibration? & Marginal row-level coverage under a held-out regime is restored by a max/similarity fallback radius (within-regime calibration is the unexercised positive control in the held-out-regime L0 design); trajectory-level $G$ or regime-segment $G_r$ coverage is recoverable only with evaluation-regime calibration data; branch-conditional coverage is recoverable with calibration-side branch labels via branch-union at a width cost. & \S5.2; \S\S7.1--7.4; Table~\ref{tab:4} and Appendix Tables B5, B7 \\
Q3. What is impossible? & For branch-blind, nondegenerate score rules, no distribution-free branch-conditional guarantee from pre-decision information; within the non-Markov freeze class without I5/I5$'$, there exist observationally equivalent models with different values when the bridge reward is not determined by endpoints (Theorem 4 is an existence result inside that class); no coverage-gain-per-unit-width lower bound. & \S5, Theorems 3--4, Propositions 2A--2C; \S\S7.3--7.6 \\
\bottomrule
\end{tabular}
\vspace{0.5em}
\end{table}

\begin{table}[!ht]
\centering
\scriptsize
\caption{Real-path usage check in the liquid-regime case}
\label{tab:7}
\par\noindent\begin{minipage}{\linewidth}\raggedright{\scriptsize Sample: 33 A-layer symbol-events over 16 independent freeze episodes (14 weekend + 2 holiday), 2026-07-01T00:00Z/2026-10-20T00:00Z. Post-hoc regime-conditional case, not a pre-registered sample.}\end{minipage}\par

\par\noindent\begin{minipage}{\linewidth}\raggedright{\scriptsize A. Observation-map parameters}\end{minipage}\par

\begin{adjustbox}{max width=\linewidth}
\begin{tabular}{@{}lrrrrr@{}}
\toprule
Symbol & Events & sigma\_open & sigma\_closed & open\_sd & stale\_phi \\
\midrule
SPYon & 15 & 0.00205 & 0.00174 & 0.00405 & 0.908 \\
NVDAon & 3 & 0.00477 & 0.00152 & 0.00320 & 0.920 \\
TSLAon & 15 & 0.00721 & 0.00337 & 0.00588 & 0.878 \\
\bottomrule
\end{tabular}
\end{adjustbox}
\vspace{0.5em}
\par\noindent\begin{minipage}{\linewidth}\raggedright{\scriptsize B. Event-level conformal coverage (alpha = 0.1, leave-one-episode-out, raw score)}\end{minipage}\par

\begin{adjustbox}{max width=\linewidth}
\begin{tabular}{@{}lrrrrrr@{}}
\toprule
Target & Identified cov. & Identified width & Identified MAE & Naive cov. & Naive width & Naive MAE \\
\midrule
Freeze & 0.909 & 0.00033 & 0.00006 & 0.909 & 0.00191 & 0.00045 \\
Post-open OPEN & 0.909 & 0.00430 & 0.00074 & 0.909 & 0.00744 & 0.00105 \\
Post-open 2 h & 0.909 & 0.00236 & 0.00032 & 0.909 & 0.00252 & 0.00032 \\
Event window & 0.909 & 0.00402 & 0.00105 & 0.909 & 0.00414 & 0.00095 \\
Total main & 0.909 & 0.00535 & 0.00122 & 0.909 & 0.00559 & 0.00104 \\
\bottomrule
\end{tabular}
\end{adjustbox}
\vspace{0.5em}
\par\noindent\begin{minipage}{\linewidth}\raggedright{\scriptsize Coverage is 30/33 in every cell; the cluster-bootstrap 95\% CI is [0.81, 1.00] for the total and event windows. Widths and MAE are in units of $V_0$.}\end{minipage}\par

\par\noindent\begin{minipage}{\linewidth}\raggedright{\scriptsize C. Modeled replay reward decomposition (Wide policy, mean truth, fraction of V0)}\end{minipage}\par

\begin{adjustbox}{max width=\linewidth}
\begin{tabular}{@{}lrrrr@{}}
\toprule
Window & Fee & IL & Cost & Net \\
\midrule
Freeze & +0.000056 & -0.000019 & +0.000000 & +0.000038 \\
Post-open 2 h & +0.000003 & -0.000107 & +0.000220 & -0.000323 \\
Event window & +0.000067 & -0.000216 & +0.000251 & -0.000399 \\
Total main & +0.000116 & -0.000653 & +0.001398 & -0.001934 \\
\bottomrule
\end{tabular}
\end{adjustbox}
\vspace{0.5em}
\par\noindent\begin{minipage}{\linewidth}\raggedright{\scriptsize P-marked model decomposition; not feeGrowthGlobal or deployed-LP validation.}\end{minipage}\par

\par\noindent\begin{minipage}{\linewidth}\raggedright{\scriptsize D. Score sensitivity (coverage / width)}\end{minipage}\par

\begin{adjustbox}{max width=\linewidth}
\begin{tabular}{@{}llrr@{}}
\toprule
Target & Score & Identified & Naive \\
\midrule
freeze & raw & 0.909 / 0.00033 & 0.909 / 0.00191 \\
freeze & studentized & 0.909 / 0.00041 & 0.909 / 0.00123 \\
freeze & cqr & 0.909 / 0.00038 & 0.909 / 0.00200 \\
total\_main & raw & 0.909 / 0.00535 & 0.909 / 0.00559 \\
total\_main & studentized & 0.909 / 0.00407 & 0.909 / 0.00388 \\
total\_main & cqr & 0.909 / 0.00447 & 0.909 / 0.00399 \\
\bottomrule
\end{tabular}
\end{adjustbox}
\vspace{0.5em}
\par\noindent\begin{minipage}{\linewidth}\raggedright{\scriptsize E. Branch and target-shift diagnostics}\end{minipage}\par

\begin{adjustbox}{max width=\linewidth}
\begin{tabular}{@{}lrr@{}}
\toprule
Diagnostic & Identified & Naive \\
\midrule
Wide branch & B=1 n=7; union n.i. & B=1 n=7; union n.i. \\
Narrow branch & B=1 n=14; union 0.970; width x1.28 & B=1 n=14; union 0.970; width x1.08 \\
Leave-one-symbol-out RAW (TSLA) & 0.267 & 0.267 \\
Leave-one-symbol-out CQR (TSLA) & 0.933 & 0.933 \\
Within-symbol (SPY/TSLA macro) & 0.933 & 0.933 \\
In-sample 2h-\textgreater{}24h target shift & 0.879 & 0.848 \\
\bottomrule
\end{tabular}
\end{adjustbox}
\vspace{0.5em}
\par\noindent\begin{minipage}{\linewidth}\raggedright{\scriptsize n.i. = not identifiable at 90\% (fewer than nine branch calibration events). The symbol-conditional row is a macro average over SPY and TSLA; NVDA has 3 events and is excluded from a 90\% within-symbol radius. Duration rows are in-sample target shifts over the same events, not out-of-sample transfer; reverse 1.000 reflects the wider 24h radius, not a symmetric guarantee.}\end{minipage}\par

\par\noindent\begin{minipage}{\linewidth}\raggedright{\scriptsize Notes - Post-hoc liquid-regime window (2026-07-01T00:00Z/2026-10-20T00:00Z); not a pre-registered sample. Results are regime-conditional and descriptive. - Truth is deterministic replay of the target policy on the observed pool-price path with a P-marked reward; it is not a deployed-policy counterfactual and not a causal logged OPE estimate. - Identified/naive are model-based OPE pipelines; the model is fit to observable P/oracle-basis data, not to latent S. - open\_sd is estimated from event-level boundary basis values, not a full Chainlink round path; treat it as a first-pass estimate. - CLOSED\_LIVE is a separate observation branch with its own sigma\_closed; its noise scale is derived from open\_sd through the closed/open volatility ratio, not estimated independently. - NVDAon uses the NVDA-USD (24/5) proxy feed and contributes only 3 of the 33 A-layer events. - Bootstrap CIs resample the 16 calendar dates, not symbol-events; they account for cross-symbol dependence but are wide. - Observation-map parameters are estimated once on all events; leave-one-episode-out folds recalibrate only the conformal radius, so this is not nested cross-validation. - Duration target-shift rows calibrate and evaluate on the same event set; they are not out-of-sample transfer results. - Thin events are retained for coverage, width, and transfer analyses; the sensitivity check splits coverage over all 33 events (16 thin/17 non-thin). No economic-value estimate is reported.}\end{minipage}\par

\end{table}

